% 自動生成（tools/mkchecklist.py p2）。各章の「ふりかえり」の項目を集めたもの（第13回）
\begingroup\footnotesize
\def\bkcb{\tikz[baseline=-0.5ex]\draw[\bkTF, line width=0.7pt, rounded corners=0.5pt] (0,-0.3em) rectangle (0.8em,0.5em);}
\begin{longtblr}{colspec={Q[c,m,wd=11mm] X[l,m] Q[c,m,wd=10mm] Q[c,m,wd=9mm] Q[c,m,wd=9mm]}, width=\linewidth, rowhead=1, colsep=3pt,
  row{1}={bg=bkVol!14, font=\sffamily\bfseries\scriptsize}, hlines={0.4pt,bkLine}, rowsep=0.7pt}
章 & できるようになったこと & ページ & 1回め & 2回め \\
{\sffamily\bfseries\color{\bkTT}熱-0} & 発電所に入る燃料の熱・出る電気・捨てる熱を、熱のバランス式と効率で数えられる。 & \pageref{ph:熱-0:sec2} & \bkcb & \bkcb \\
 & 何を無視したかを書いてから、捨てる熱を運ぶ海の水の量を、7つの段で求められる。 & \pageref{ph:熱-0:ex3} & \bkcb & \bkcb \\
 & 身近な「なぜ？」の答えが、この巻のどの章にあるかを探せる。 & \pageref{ph:熱-0:sec6} & \bkcb & \bkcb \\
{\sffamily\bfseries\color{\bkTT}熱-1} & 氷がとけている間、温度が変わらないわけを、潜熱で説明できる。 & \pageref{ph:熱-1:sec5} & \bkcb & \bkcb \\
 & $Q=mc\,\Delta T$ と熱のバランス式で、温まり方と混ぜたときの温度を求められる。 & \pageref{ph:熱-1:sec4} & \bkcb & \bkcb \\
 & 加熱曲線を読み、熱膨張の大きさを見積もれる。 & \pageref{ph:熱-1:sec6} & \bkcb & \bkcb \\
 & 燃料が燃えるときの空気の量と発熱量を、原子のバランスで求められる。 & \pageref{ph:熱-1:sec8} & \bkcb & \bkcb \\
{\sffamily\bfseries\color{\bkTT}熱-2} & 熱と仕事が同じエネルギーであることを、実験から説明できる。 & \pageref{ph:熱-2:sec1} & \bkcb & \bkcb \\
 & 第一法則 $\Delta U=Q+W$ で、エネルギーの出入りを数えられる。 & \pageref{ph:熱-2:sec4} & \bkcb & \bkcb \\
 & $p$-$V$ 図の面積から、気体がする仕事を求められる。 & \pageref{ph:熱-2:sec3} & \bkcb & \bkcb \\
 & 状態量と、道すじで変わる量を、全微分で見分けられる。 & \pageref{ph:熱-2:sec7} & \bkcb & \bkcb \\
{\sffamily\bfseries\color{\bkTT}熱-3} & ボイル・シャルルの法則と状態方程式で、気体の圧力・体積・温度を求められる。 & \pageref{ph:熱-3:sec2} & \bkcb & \bkcb \\
 & 定積・定圧・等温・断熱の4つの変化で、熱と仕事と内部エネルギーの出入りを数えられる。 & \pageref{ph:熱-3:sec4} & \bkcb & \bkcb \\
 & 断熱変化・ポリトロープ変化の式で、空気入れ・エンジン・圧縮機・大気の温度と圧力を計算できる。 & \pageref{ph:熱-3:sec6} & \bkcb & \bkcb \\
{\sffamily\bfseries\color{\bkTT}熱-4} & 定常流動のエネルギーの式で、タービン・圧縮機・ノズル・ドライヤーのエネルギーを計算できる。 & \pageref{ph:熱-4:sec3} & \bkcb & \bkcb \\
 & 音の速さ・マッハ数・よどみ点の温度を計算できる。 & \pageref{ph:熱-4:sec6} & \bkcb & \bkcb \\
 & 先細末広ノズルの、のどの状態と出口の面積を求められる。 & \pageref{ph:熱-4:sec8} & \bkcb & \bkcb \\
{\sffamily\bfseries\color{\bkTT}熱-5} & 熱機関の熱効率と、冷凍機・ヒートポンプの成績係数を、熱と仕事の矢印から計算できる。 & \pageref{ph:熱-5:sec2} & \bkcb & \bkcb \\
 & 第二法則の2つの言い方を説明し、カルノーの効率 $1-T_{\mathrm L}/T_{\mathrm H}$ を求められる。 & \pageref{ph:熱-5:sec3} & \bkcb & \bkcb \\
 & カルノーの定理から、熱機関と冷凍機の上限と、熱力学温度の考えを導ける。 & \pageref{ph:熱-5:sec5} & \bkcb & \bkcb \\
{\sffamily\bfseries\color{\bkTT}熱-6} & 熱÷温度 $\Delta S=Q/T$ を数えて、もどらない向きを言える。 & \pageref{ph:熱-6:sec2} & \bkcb & \bkcb \\
 & 水・氷・気体の $\Delta S$ を計算し、$T$-$s$ 線図の面積を読める。 & \pageref{ph:熱-6:sec7} & \bkcb & \bkcb \\
 & 生まれるエントロピーとエクセルギーで、機械の損失を数えられる。 & \pageref{ph:熱-6:sec9} & \bkcb & \bkcb \\
{\sffamily\bfseries\color{\bkTT}熱-7} & 4ストロークのエンジンの動きを $p$-$V$ 図で説明し、断熱圧縮のあとの温度を求められる。 & \pageref{ph:熱-7:sec1} & \bkcb & \bkcb \\
 & オットーサイクルとディーゼルサイクルの理論熱効率を、圧縮比から計算できる。 & \pageref{ph:熱-7:sec3} & \bkcb & \bkcb \\
 & ガスタービンの仕事比・再生・コンバインドサイクルの効率と、エンジンの出力を計算できる。 & \pageref{ph:熱-7:sec6} & \bkcb & \bkcb \\
{\sffamily\bfseries\color{\bkTT}熱-8} & 水のわく温度が、まわりの圧力（飽和蒸気圧）で決まることを説明できる。 & \pageref{ph:熱-8:sec1} & \bkcb & \bkcb \\
 & 乾き度と蒸気表（直線補間）で、湿り蒸気・過熱蒸気の比エンタルピーを求められる。 & \pageref{ph:熱-8:sec2} & \bkcb & \bkcb \\
 & ランキンサイクルの効率を計算し、再熱・再生の働きを説明できる。 & \pageref{ph:熱-8:sec6} & \bkcb & \bkcb \\
{\sffamily\bfseries\color{\bkTT}熱-9} & 湿度と露点を使って、コップや窓に水がつくわけを説明できる。 & \pageref{ph:熱-9:sec6} & \bkcb & \bkcb \\
 & 絶対湿度・比エンタルピー・空気線図で、エアコンが空気から取る水と熱を計算できる。 & \pageref{ph:熱-9:sec5} & \bkcb & \bkcb \\
 & 冷凍サイクルの成績係数（COP）を、冷媒の表から計算できる。 & \pageref{ph:熱-9:sec2} & \bkcb & \bkcb \\
{\sffamily\bfseries\color{\bkTT}熱-10} & 蒸気表とファンデルワールスの式で、わく温度と実在気体の圧力を求められる。 & \pageref{ph:熱-10:sec8} & \bkcb & \bkcb \\
 & $\ln p$ と $1/T$ のグラフから蒸発熱を読み、わく温度を見積もれる。 & \pageref{ph:熱-10:sec3} & \bkcb & \bkcb \\
 & 自由エネルギーとマクスウェルの関係式から、相平衡の式を導ける。 & \pageref{ph:熱-10:sec7} & \bkcb & \bkcb \\
{\sffamily\bfseries\color{\bkTT}熱-11} & 熱の伝わり方（伝導・対流・放射）を、身近な例で見分けられる。 & \pageref{ph:熱-11:sec1} & \bkcb & \bkcb \\
 & フーリエの法則で、平らな壁を通りぬける熱の流れを求められる。 & \pageref{ph:熱-11:sec2} & \bkcb & \bkcb \\
 & 熱抵抗の直列と熱貫流率で、窓や壁から逃げる熱と、表面の温度を求められる。 & \pageref{ph:熱-11:sec4} & \bkcb & \bkcb \\
 & 壁を小さな箱に分けた表計算で、壁の中の温度の変わり方を追える。 & \pageref{ph:熱-11:sec5} & \bkcb & \bkcb \\
 & 管の保温材の熱抵抗と、ビオ数・時定数を使って、冷え方を見積もれる。 & \pageref{ph:熱-11:sec8} & \bkcb & \bkcb \\
 & 熱方程式を導き、板の冷え方や地中の温度を、変数分離の答えで読める。 & \pageref{ph:熱-11:sec9} & \bkcb & \bkcb \\
{\sffamily\bfseries\color{\bkTT}熱-12} & ニュートンの冷却則と熱伝達率の目安で、物から流れへ逃げる熱を見積もれる。 & \pageref{ph:熱-12:sec2} & \bkcb & \bkcb \\
 & 冷めていく温度の測定から、表計算で熱伝達率を求められる。 & \pageref{ph:熱-12:sec3} & \bkcb & \bkcb \\
 & フィンを小さな箱に分けた表計算で、フィンの温度と効きめを求められる。 & \pageref{ph:熱-12:sec4} & \bkcb & \bkcb \\
 & 無次元数と相関式を選んで、強制対流・自然対流の熱伝達率を求められる。 & \pageref{ph:熱-12:sec6} & \bkcb & \bkcb \\
 & 沸騰曲線と限界熱流束、凝縮の熱伝達を説明し、大きさを見積もれる。 & \pageref{ph:熱-12:sec9} & \bkcb & \bkcb \\
{\sffamily\bfseries\color{\bkTT}熱-13} & 物が出す放射の強さ $\sigma T^4$ と、いちばん強い波長を計算できる。 & \pageref{ph:熱-13:sec2} & \bkcb & \bkcb \\
 & 放射率と形態係数を使って、2つの面の間で放射で移る熱を計算できる。 & \pageref{ph:熱-13:sec6} & \bkcb & \bkcb \\
 & 対数平均温度差と温度効率で、熱交換器の熱の量と大きさを求められる。 & \pageref{ph:熱-13:sec8} & \bkcb & \bkcb \\
{\sffamily\bfseries\color{\bkTT}熱-14} & 分子が壁にぶつかる力積から、気体の圧力と温度の意味を説明できる。 & \pageref{ph:熱-14:sec2} & \bkcb & \bkcb \\
 & 気体の内部エネルギーと比熱を、分子の運動から計算できる。 & \pageref{ph:熱-14:sec4} & \bkcb & \bkcb \\
 & 速さの分布と平均自由行程から、粘度・熱伝導率・拡散を見積もれる。 & \pageref{ph:熱-14:sec5} & \bkcb & \bkcb \\
{\sffamily\bfseries\color{\bkTT}熱-15} & 並び方の数を数えて、変化がなぜ一方の向きにだけ進むのかを説明できる。 & \pageref{ph:熱-15:sec2} & \bkcb & \bkcb \\
 & $S=k_{\mathrm B}\ln\Omega$ で、並び方の数からエントロピーを計算できる。 & \pageref{ph:熱-15:sec4} & \bkcb & \bkcb \\
 & ボルツマン分布と分配関数から、2準位の系や固体の比熱を導ける。 & \pageref{ph:熱-15:sec7} & \bkcb & \bkcb \\
{\sffamily\bfseries\color{\bkTT}熱-16} & \qty{1}{kWh}の電気に要る燃料の熱と、燃やして出る二酸化炭素の量を、効率と原子の数のバランスから見積もれる。 & \pageref{ph:熱-16:sec2} & \bkcb & \bkcb \\
 & 発電や給湯の方法を、効率・温度・設備利用率・一生で出る二酸化炭素など、いくつもの物差しでくらべられる。 & \pageref{ph:熱-16:sec9} & \bkcb & \bkcb \\
 & ヒートポンプや燃料電池が得をするわけを、エクセルギー（第二法則）で説明できる。 & \pageref{ph:熱-16:sec6} & \bkcb & \bkcb \\
\end{longtblr}
\endgroup
